\documentclass[11pt]{article}
\usepackage[margin=1in]{geometry}
\usepackage{booktabs}
\usepackage{amsmath,amssymb}
\usepackage{array}
\usepackage{xcolor}
\usepackage{hyperref}
\usepackage{caption}
\usepackage{subcaption}
\usepackage{graphicx}
\usepackage{multirow}
\graphicspath{{output/}{output_scaling/}{exp_onesided/output/}{exp_onesided/output_scaling/}}

\title{One-Sided Quantile Estimation: CKME vs.\ Linear Quantile Regression\\[0.5em]
	\large Experiment Summary and Analysis}
\author{}
\date{\today}

\begin{document}
	\maketitle
	
	% ===========================================================================
	\section{Methods}
	% ===========================================================================
	
	\subsection{Quantile estimators}
	
	\begin{description}
		\item[CKME (CDF-first).] We first fit a conditional kernel mean embedding model
		$\hat{F}(y \mid x)$ on the training data, and then invert the estimated CDF
		on a fixed $t$-grid to obtain
		\[
		\hat{q}_\tau(x) = \inf\{y : \hat{F}(y \mid x) \ge \tau\}.
		\]
		\item[QR (quantile-first).] We fit linear quantile regression at level $\tau$
		directly using degree-3 polynomial features
		\[
		\phi(x) = (1, x, x^2, x^3)^\top
		\]
		through \texttt{statsmodels.QuantReg}.
	\end{description}
	
	\subsection{One-sided bounds and target coverage}
	
	For $\tau \in \{0.05, 0.95\}$, we consider the lower and upper one-sided quantile bounds:
	\begin{align*}
		\text{Lower bound:}\quad & L(x) = \hat{q}_{0.05}(x),
		\qquad \text{target } P(Y \ge L(X)) \approx 0.95, \\
		\text{Upper bound:}\quad & U(x) = \hat{q}_{0.95}(x),
		\qquad \text{target } P(Y \le U(X)) \approx 0.95.
	\end{align*}
	
	Thus, the lower-tail problem focuses on estimating the $0.05$ quantile accurately,
	while the upper-tail problem focuses on the $0.95$ quantile.
	
	\subsection{Evaluation metrics}
	
	\begin{description}
		\item[Empirical coverage.] For the lower bound, we report
		$\hat{P}(Y \ge \hat{q}_{0.05}(X))$; for the upper bound, we report
		$\hat{P}(Y \le \hat{q}_{0.95}(X))$, both evaluated on the test set.
		\item[Pinball loss.] For a quantile estimate $q$ at level $\tau$, the pinball loss is
		\[
		\rho_\tau(y, q)
		= (y-q)\bigl[\tau \mathbf{1}_{y \ge q} + (\tau-1)\mathbf{1}_{y<q}\bigr],
		\]
		averaged over test points.
		\item[Mean absolute quantile error.]
		\[
		\overline{|e|}
		= \frac{1}{n_{\mathrm{test}}}
		\sum_{i=1}^{n_{\mathrm{test}}}
		\bigl|\hat{q}_\tau(x_i) - q_\tau^*(x_i)\bigr|,
		\]
		where $q_\tau^*(x)$ denotes the Monte Carlo oracle quantile.
		\item[Sup absolute quantile error.]
		\[
		\sup_i \bigl|\hat{q}_\tau(x_i)-q_\tau^*(x_i)\bigr|.
		\]
	\end{description}
	
	% ===========================================================================
	\section{Experiment 1: CKME vs.\ QR Comparison}
	% ===========================================================================
	
	\subsection{Settings}
	
	\begin{center}
		\begin{tabular}{ll}
			\toprule
			Parameter & Value \\
			\midrule
			Training sites & $n_{\mathrm{train}} = 1000$ \\
			Replications per site & $r_{\mathrm{train}} = 1$ (i.i.d.) \\
			Test points & $n_{\mathrm{test}} = 1000$ (uniform grid) \\
			$t$-grid size & 500 \\
			QR polynomial degree & 3 \\
			Monte Carlo samples for true quantiles & 10{,}000 per $x$ \\
			Macroreplications & 5 \\
			\bottomrule
		\end{tabular}
	\end{center}
	
	\subsection{Simulators}
	
	\begin{center}
		\begin{tabular}{lll}
			\toprule
			Name & True function & Noise \\
			\midrule
			\texttt{exp1}
			& $f(x)=1.5x^2/(1-x)$,\quad $x\in[0.1,0.9]$
			& Heteroscedastic Gaussian \\
			\texttt{exp2}
			& $f(x)=x+\sin(\pi x)$,\quad $x\in[0,2\pi]$
			& Heteroscedastic Gaussian \\
			\texttt{nongauss\_B2L}
			& $f(x)=e^{x/10}\sin(x)$,\quad $x\in[0,2\pi]$
			& Centered Gamma ($k=2$, strongly skewed) \\
			\bottomrule
		\end{tabular}
	\end{center}
	
	For CKME, hyperparameters are loaded from \texttt{pretrained\_params.json}
	using prior CV tuning:
	\[
	\texttt{exp1}: (\ell_x,\lambda,h)=(0.1,0.001,0.1),\qquad
	\texttt{exp2}: (0.5,0.001,0.1),\qquad
	\texttt{nongauss\_B2L}: (0.5,0.001,0.1).
	\]
	
	\subsection{Mean absolute quantile error (averaged over 5 macroreplications)}
	
	\begin{table}[h]
		\centering
		\caption{Mean and sup absolute quantile error over the test grid.
			Boldface indicates the smaller mean error within each simulator and tail level.}
		\begin{tabular}{llcccc}
			\toprule
			& & \multicolumn{2}{c}{$\tau=0.05$ (lower)} &
			\multicolumn{2}{c}{$\tau=0.95$ (upper)} \\
			\cmidrule(lr){3-4}\cmidrule(lr){5-6}
			Simulator & Method & Mean $|\hat{q}-q^*|$ & Sup $|\hat{q}-q^*|$ &
			Mean $|\hat{q}-q^*|$ & Sup $|\hat{q}-q^*|$ \\
			\midrule
			\multirow{2}{*}{\texttt{exp1}}
			& CKME & $0.145$ & $0.93$  & $\mathbf{0.653}$ & $30.00$ \\
			& QR   & $\mathbf{0.083}$ & $2.19$ & $0.667$ & $\mathbf{7.70}$ \\
			\midrule
			\multirow{2}{*}{\texttt{exp2}}
			& CKME & $0.148$ & $0.46$  & $0.278$ & $4.16$ \\
			& QR   & $\mathbf{0.104}$ & $0.89$ & $\mathbf{0.114}$ & $\mathbf{0.75}$ \\
			\midrule
			\multirow{2}{*}{\texttt{nongauss\_B2L}}
			& CKME & $0.150$ & $0.33$  & $0.233$ & $4.35$ \\
			& QR   & $\mathbf{0.099}$ & $0.87$ & $\mathbf{0.103}$ & $\mathbf{0.80}$ \\
			\bottomrule
		\end{tabular}
	\end{table}
	
	\begin{figure}[!htbp]
		\centering
		\includegraphics[width=0.85\textwidth]{onesided_mean_error_boxplot.png}
		\caption{Mean absolute quantile error $\overline{|\hat{q}_\tau-q_\tau^*|}$ across 5 macroreplications.
			Left: $\tau=0.05$ (lower bound). Right: $\tau=0.95$ (upper bound).}
		\label{fig:mean_error}
	\end{figure}
	
	\begin{figure}[!htbp]
		\centering
		\includegraphics[width=0.85\textwidth]{onesided_sup_error_boxplot.png}
		\caption{Sup absolute quantile error $\sup_i |\hat{q}_\tau(x_i)-q_\tau^*(x_i)|$ across 5 macroreplications.
			The particularly large CKME sup error at $\tau=0.95$ for \texttt{exp1} (around $30$)
			indicates a severe local failure at a small set of $x$ values.}
		\label{fig:sup_error}
	\end{figure}
	
	\paragraph{Key observations.}
	First, QR attains smaller mean quantile error than CKME for all three simulators at
	$\tau=0.05$. Second, for $\tau=0.95$, QR remains clearly better on \texttt{exp2}
	and \texttt{nongauss\_B2L}, whereas on \texttt{exp1} the two methods are similar in
	\emph{mean} error but differ sharply in \emph{worst-case} behavior: CKME has a sup
	error around $30$, compared with $7.7$ for QR. This indicates that the average
	performance of CKME at the upper tail can mask catastrophic failure at a few specific
	test locations.
	
	A second pattern is a clear tail asymmetry for CKME. Across all simulators, its sup
	error is substantially larger for $\tau=0.95$ than for $\tau=0.05$, whereas QR is
	much more balanced across the two tails. This suggests that the main weakness of CKME
	in this setting is concentrated in upper-tail quantile estimation rather than being a
	uniform issue over the full distribution.
	
	% ===========================================================================
	\section{Experiment 2: CKME Scaling Study}
	% ===========================================================================
	
	\subsection{Settings}
	
	\begin{center}
		\begin{tabular}{ll}
			\toprule
			Parameter & Value \\
			\midrule
			Simulators & \texttt{exp1}, \texttt{exp2}, \texttt{nongauss\_A1L} \\
			$n_{\mathrm{train}}$ & $\{100,200,500,1000,2000\}$ \\
			$r_{\mathrm{train}}$ & $1$ (one observation per site, i.i.d.) \\
			$n_{\mathrm{test}}$ & 500 (uniform grid) \\
			$t$-grid size & 500 \\
			Monte Carlo samples for true quantiles & 10{,}000 per $x$ \\
			Macroreplications & 10 \\
			\bottomrule
		\end{tabular}
	\end{center}
	
	Here \texttt{nongauss\_A1L} uses
	\[
	f(x)=e^{x/10}\sin(x), \qquad x\in[0,2\pi],
	\]
	with Student-$t$ noise ($\nu=3$), producing a heavy-tailed conditional distribution.
	
	\subsection{Mean absolute quantile error vs.\ $n_{\mathrm{train}}$}
	
	The values below are $\overline{|e|}$ averaged over 10 macroreplications.
	
	\subsubsection{$\tau = 0.05$ (lower bound)}
	\begin{center}
		\begin{tabular}{lccccc}
			\toprule
			Simulator & $n=100$ & $n=200$ & $n=500$ & $n=1000$ & $n=2000$ \\
			\midrule
			\texttt{exp1}          & 0.401 & 0.341 & 0.253 & 0.217 & 0.182 \\
			\texttt{exp2}          & 0.443 & 0.362 & 0.267 & 0.220 & 0.175 \\
			\texttt{nongauss\_A1L} & 0.331 & 0.294 & 0.217 & 0.171 & 0.148 \\
			\bottomrule
		\end{tabular}
	\end{center}
	
	\subsubsection{$\tau = 0.95$ (upper bound)}
	\begin{center}
		\begin{tabular}{lccccc}
			\toprule
			Simulator & $n=100$ & $n=200$ & $n=500$ & $n=1000$ & $n=2000$ \\
			\midrule
			\texttt{exp1}          & 0.730 & 0.686 & 0.654 & 0.655 & 0.663 \\
			\texttt{exp2}          & 0.428 & 0.370 & 0.304 & 0.273 & 0.264 \\
			\texttt{nongauss\_A1L} & 0.326 & 0.332 & 0.292 & 0.280 & 0.292 \\
			\bottomrule
		\end{tabular}
	\end{center}
	
	\subsection{Sup error vs.\ $n_{\mathrm{train}}$ for $\tau=0.95$}
	
	\begin{center}
		\begin{tabular}{lccccc}
			\toprule
			Simulator & $n=100$ & $n=200$ & $n=500$ & $n=1000$ & $n=2000$ \\
			\midrule
			\texttt{exp1}          & 13.8 & 15.9 & 20.5 & 23.5 & 27.2 \\
			\texttt{exp2}          & 2.09 & 1.94 & 2.10 & 2.34 & 3.09 \\
			\texttt{nongauss\_A1L} & 1.55 & 2.19 & 3.20 & 4.57 & 5.94 \\
			\bottomrule
		\end{tabular}
	\end{center}
	
	\begin{figure}[htbp]
		\centering
		\includegraphics[width=\textwidth]{onesided_tailprob_fixed_t_agg.png}
		\caption{Estimated tail probability $\hat{F}(t \mid x)$ at a fixed threshold $t$,
			aggregated across macroreplications. Deviations from the oracle curve indicate
			systematic bias in CKME's estimated upper tail.}
		\label{fig:tailprob_agg}
	\end{figure}
	
	\subsection{Diagnostic summary}
	
	Three diagnostics were run using \texttt{diag\_ckme\_scaling.py}.
	
	\paragraph{Diag~1 (worst-case location $x^*$).}
	For $\tau=0.05$, the worst-case location
	\[
	x^*=\arg\max_x |\hat{q}_\tau(x)-q_\tau^*(x)|
	\]
	is spread across the domain and gradually contracts as $n$ increases.
	For $\tau=0.95$, however, $x^*$ consistently concentrates near the right boundary
	$x=2\pi$, indicating that the upper-tail failure is not diffuse but localized in a
	specific region.
	
	\begin{figure}[htbp]
		\centering
		\includegraphics[width=\textwidth]{diag_worst_x.png}
		\caption{Diag~1: distribution of worst-case locations
			$x^*=\arg\max_x |\hat{q}_\tau(x)-q_\tau^*(x)|$ across macroreplications, for each
			$n_{\mathrm{train}}$. Top: $\tau=0.05$. Bottom: $\tau=0.95$.}
		\label{fig:diag_worst_x}
	\end{figure}
	
	\paragraph{Diag~4 (error percentiles).}
	For $\tau=0.05$, the mean, 95th percentile, and sup errors all decrease with
	$n_{\mathrm{train}}$, indicating standard convergence behavior.
	
	For $\tau=0.95$, the behavior is more heterogeneous:
	\begin{itemize}
		\item For \texttt{exp2} and \texttt{nongauss\_A1L}, the 95th percentile error stays
		close to the mean error, suggesting that most test points improve with larger
		$n$, and that the growing sup error is driven by a small number of outlying
		$x$ values.
		\item For \texttt{exp1}, the pattern is much more anomalous: the mean error remains
		nearly flat at about $0.66$, while the sup error increases steadily from $13.8$
		at $n=100$ to $27.2$ at $n=2000$. In other words, adding more data does not
		improve upper-tail accuracy and in fact worsens the worst-case error.
	\end{itemize}
	
	\begin{figure}[htbp]
		\centering
		\includegraphics[width=\textwidth]{diag_q95err.png}
		\caption{Diag~4: mean, 95th percentile, and sup error of
			$|\hat{q}_{0.95}(x)-q_{0.95}^*(x)|$ versus $n_{\mathrm{train}}$.
			For \texttt{exp1}, the increasing sup error confirms anomalous upper-tail divergence.}
		\label{fig:diag_q95err}
	\end{figure}
	
	% ===========================================================================
	\section{Discussion}
	% ===========================================================================
	
	\subsection{Upper-tail failure of CKME}
	
	The main finding is a pronounced asymmetry between the lower and upper tails for CKME.
	
	For the lower tail ($\tau=0.05$), CKME shows regular convergence in mean error as
	$n_{\mathrm{train}}$ increases. Although it is somewhat conservative from a coverage
	perspective, its estimation behavior is broadly stable.
	
	For the upper tail ($\tau=0.95$), the picture is very different. On \texttt{exp1},
	the mean error does not meaningfully decrease with $n$, and the sup error increases
	substantially. On the other simulators, the upper-tail mean error does improve, but
	much more slowly than in the lower tail, and the worst-case error still grows.
	
	This pattern suggests that the CKME difficulty is not simply due to finite-sample
	noise. Rather, the issue appears structural and is concentrated in the upper tail,
	especially near certain boundary regions.
	
	\subsection{A plausible explanation}
	
	The diagnostics suggest that the problem is not caused by an insufficiently wide
	$t$-grid. Instead, a more plausible explanation is that the smooth-indicator CKME
	approximation allocates its fitting effort too uniformly over $t$ and therefore does
	not represent extreme upper-tail behavior accurately enough.
	
	In particular, if $\hat{F}(y \mid x)$ grows too slowly near high probability levels,
	then the inversion
	\[
	\hat{q}_{0.95}(x)=\inf\{y:\hat{F}(y\mid x)\ge 0.95\}
	\]
	will be shifted upward relative to the oracle quantile. This leads to overly large
	upper bounds, inflated one-sided coverage, and large quantile estimation error.
	With more data, the central part of the conditional CDF may improve, while the tail
	region remains poorly captured or even becomes relatively more distorted under the
	global fitting criterion.
	
	\subsection{Why QR performs better here}
	
	Linear polynomial QR directly optimizes the pinball loss at the target quantile level,
	so it is tuned specifically for the quantity being evaluated. By contrast, CKME is
	estimated through a full-CDF objective, which spreads fitting effort across the
	entire distribution rather than prioritizing a single tail quantile.
	
	For the Gaussian examples \texttt{exp1} and \texttt{exp2}, the degree-3 QR model is
	also reasonably well specified or only mildly misspecified. In that setting, there is
	little reason to expect a full-distribution method to outperform a model that targets
	the quantile directly. The non-Gaussian example \texttt{nongauss\_B2L} still favors
	QR, which suggests that, at least in the current implementation, CKME's flexibility
	does not translate into better plug-in one-sided quantile estimation.
	
	\subsection{Implications for the paper}
	
	At the plug-in estimation level, without conformal calibration, the current evidence
	supports the following conclusions.
	
	First, QR generally provides more accurate one-sided quantile estimates and lower
	tail-specific loss. Second, CKME tends to be more conservative, especially for upper
	quantiles, because its estimated conditional distribution can be overly spread in the
	relevant tail region. Third, the main value of CKME is therefore unlikely to be raw
	plug-in quantile accuracy. Its real advantage lies in representing the full
	conditional law, which is useful for downstream tasks such as conformal calibration,
	distributional scoring, and adaptive data collection in the two-stage pipeline.
	
	This suggests an important framing choice for the paper: CKME should not be presented
	as a universally better plug-in quantile estimator than QR. Instead, it is better
	motivated as a full-distribution model whose strength appears after conformalization
	or within broader distributional decision pipelines.
	
\end{document}